\subsubsection{NR and NS: Replacing the Evaluator, Not Just Its Queue}
\label{sec:tw-variants}

The original \code{NR} and \code{NS} controllers in
\filepath{tw-wash1.abclr2} explore different implementations of the
Time Warp evaluator. They are not additional names for the priority
and FIFO policies already compared by \code{starter}.
The latter two replace the \emph{queue} of one schedulable evaluator;
NR and NS install entirely new evaluator objects.

\begin{longtable}{L{0.15\textwidth}L{0.36\textwidth}L{0.39\textwidth}}
\textbf{Variant} & \textbf{Evaluator's metaobject} & \textbf{Meaning of the change}\\\hline
\endhead
Standard & Created by \code{TW-meta-eval-gen}. & Its evaluation-request
queue can be selected as Lowest-LVT-first or FIFO.\\
NR & Explicit \code{non-reifying-meta}. & Evaluation requests use direct,
non-reifying handlers instead of the custom scheduling metaobject.\\
NS & Default ordinary metaobject. & Removes the custom Time Warp
evaluation-request scheduler, but keeps ordinary object dispatch.\\
\end{longtable}

All three evaluator definitions retain a local virtual-time cache and
the extended \code{Outputs}/\code{LVT} context. All are used with
\code{TW-meta-gen} for simulation objects, so selecting NR does not
remove those objects' rollback histories. ``No scheduling'' in NS
does not mean that the runtime stops scheduling processes or ignores
simulation event times. It refers to the omitted customized evaluator
scheduling layer; ordinary queues and the undoable event queue remain.

The group selectors show the replacement directly:

\wholelisting[firstline=22,lastline=31]{examples/time-warp/tw-group.abclr2}

Each creates a fresh evaluator and assigns the group's \code{evaluator}
binding. Unlike \code{:change-scheduler}, these handlers do not transfer
the old evaluator's pending request queue. Nor do they rewrite the
\code{evaluator} binding held by an existing simulation metaobject.
The checks create an object before switching and another afterwards:
the former retains the old evaluator, the latter captures the new one.
This is why the original NR/NS controller switches before
\code{starter-1} creates its attendant and washers. A group's configuration
is a creation policy, not necessarily a live broadcast to all its members.
Use separate fresh sessions for these experiments; restoring the group's
standard evaluator does not retrofit older simulation objects.

\paragraph{Read the controller's result protocol carefully.}
Here is NR's original controller:

\wholelisting[firstline=127,lastline=142]{examples/time-warp/tw-wash1.abclr2}

It selects the new evaluator, resets the rollback monitor, and waits on
the familiar \code{starter-1} simulation. If that request completes,
it prints the result, reads the rollback counter, and replies
\code{[:fifo [:rollback count]]}. NS follows the same structure but sends
\code{:no-schedule-eval}. Both reuse the historical \code{:fifo} result
tag, even though neither installs \code{normal-queue-gen} here. The tag
therefore cannot identify the selected evaluator architecture.
The expression-counter reset/query and \code{:exps} result fields are
commented out. The custom counter belongs to \code{TW-meta-eval}; an
ordinary or non-reifying evaluator metaobject does not acquire its
\code{:count} handler merely because the group has that query method.

The original entry points, after loading \code{tw-load} in a scratch
directory, are:

\begin{codeblock}
\begin{lstlisting}
; Use a fresh session for each variant.
[NR <== [:start [3 4 6] 3]]
; In another fresh session:
[NS <== [:start [3 4 6] 3]]
\end{lstlisting}
\end{codeblock}

Both complete this small case. Use the returned controller result to
identify completion; printed car assignments alone are not a join.

\paragraph{Follow the evaluator's representations.}
The translator represents a simple literal such as \code{:fin} as
\code{[:self-evaluate function]}, where the zero-argument function
returns the literal. Each evaluator invokes that function with
\code{funcall}; a function-valued payload is returned as data.
The standard evaluator shows the protocol directly:

\wholelisting[firstline=43,lastline=44]{examples/time-warp/tw-eval.abclr2}

The \code{[:script scriptObject]} case delegates to
\code{call-script-body}, which accepts internal scripts, object-level
\code{:exec}, or compiled functions. Trace the representation before
deciding which case handles a request.

\paragraph{What the variant checks establish.}
Run \code{make -C tutorial timewarp-variants-check}. Fifteen probes per
variant check a fresh evaluator identity, old/new objects' captured
evaluators, the literal's returned representation and payload,
a synthetic five-argument compiled body that replies 42, and
\code{LVT} lookup at 7. The compiled-body probe exercises
\code{:compiled}, whereas the literal probe exercises
\code{:self-evaluate}; one working case does not certify the other.
Further checks cover numbers, nil, symbols, lists, function-valued payloads,
single invocation of an initializer, and an application script returning 42.

Two car-wash sessions add four checks each: controller result shape,
historical \code{:fifo} tag, and a nonnegative rollback count. Transcripts
also require the \code{:fin} result and all three expected car assignments.
The 53 checks cover these evaluator and application protocols.
Inputs and transcripts stay in \filepath{build/timewarp-variants/}.
Use small runs to understand the control flow before exploring
additional cancellation orders or larger simulations.

\paragraph{Check your understanding.}
Does NR eliminate application-state rollback? No; simulation objects
still use the Time Warp metaobject generator. Does NS mean no queues?
No; it omits a particular scheduling layer. Does a returned \code{:fifo}
tag prove FIFO evaluator scheduling? No; the old controllers reuse that
tag. Can a new group evaluator fix an already created object's evaluator?
Not through these selectors: existing metaobjects retain their captured evaluator.
